% !TEX root = ../main.tex
\begin{ouverture}[gazparfait]{12}{Description d'un fluide au repos}{Partie 2 · Mouvement et interactions}
{Masse volumique\enspace·\enspace Pression\enspace·\enspace Température\enspace·\enspace Force pressante\enspace·\enspace Loi de Mariotte\enspace·\enspace Statique des fluides}
\objectif{Relier les grandeurs macroscopiques d'un fluide (masse volumique, pression, température) au comportement microscopique de ses entités.}
\objectif{Utiliser la relation $F=P\times S$ pour déterminer une force pressante.}
\objectif{Utiliser et tester la loi de Mariotte.}
\objectif{Utiliser et tester la loi fondamentale de la statique des fluides : $P_2-P_1=\rho\,g\,(z_1-z_2)$.}
\prerequis{Masse volumique, états de la matière (classe de seconde).}
\prerequis{Forces, poids $P=m\,g$ ; conversions d'unités de surface et de volume.}
\end{ouverture}
\essentiel
\essSection{1}{Décrire un fluide : deux échelles}
\begin{definition}[Grandeurs macroscopiques]
Un \textbf{fluide} (liquide ou gaz) se décrit à notre échelle par trois grandeurs :
\begin{itemize}
\item la \textbf{masse volumique} $\rho=\dfrac{m}{V}$, en \unit{kg.m^{-3}} (eau : \qty{1,00e3}{kg.m^{-3}} ; air : \qty{1,2}{kg.m^{-3}}) ;
\item la \textbf{pression} $P$, en pascals (Pa), mesurée avec un manomètre ou un capteur ; $\qty{1}{hPa}=\qty{100}{Pa}$, $\qty{1}{bar}=\qty{1e5}{Pa}$ ;
la pression atmosphérique normale vaut $P_\mathrm{atm}=\qty{1013}{hPa}$ ;
\item la \textbf{température} : $T$ (en kelvins, K) $=\theta$ (en \unit{\degreeCelsius}) $+\,273{,}15$.
\end{itemize}
\end{definition}
\begin{propriete}[Interprétation microscopique]
Les entités d'un fluide (molécules, atomes) sont en \textbf{mouvement incessant et désordonné}. La \textbf{pression} résulte de
leurs chocs sur les parois ; la \textbf{température} mesure leur agitation : plus elle est élevée, plus les entités sont rapides.
Un gaz est peu dense, car ses entités sont très éloignées les unes des autres ; elles sont au contact dans un liquide.
\end{propriete}
\schema{12-micro}
\essSection{2}{Force pressante}
\begin{loi}[Force pressante]
Un fluide à la pression $P$ exerce sur une surface plane d'aire $S$ une \textbf{force pressante} perpendiculaire à la surface,
orientée du fluide vers la surface, de valeur $\boldsymbol{F=P\times S}$, avec $F$ en N, $P$ en Pa et $S$ en \unit{m^2}. L'air exerce ainsi sur une vitre de \qty{1}{m^2} une force d'environ \qty{1,0e5}{N}, le poids d'une masse de 10 tonnes : elle est
compensée par la force de l'air situé de l'autre côté.
\end{loi}
\par\penalty0
\schema{12-force-pressante}
\essSection{3}{La loi de Mariotte}
\begin{loi}[Loi de Mariotte]
Pour une quantité de gaz donnée, à \textbf{température constante}, le produit de la pression par le volume reste constant :
\[P\times V=\text{constante}\qquad\text{soit}\qquad P_1\,V_1=P_2\,V_2.\]
Cette loi décrit bien les gaz aux pressions usuelles. Exemple : l'air d'une seringue bouchée, à \qty{1013}{hPa} sous \qty{60}{mL},
comprimé à \qty{20}{mL}, atteint $\dfrac{1013\times60}{20}=\qty{3039}{hPa}$.
\end{loi}
\par\penalty0
\schema{12-mariotte}
\begin{python}
P1, V1 = 1013, 60          # hPa, mL : air de la seringue
for V in [60, 50, 40, 30, 20]:
    P = P1 * V1 / V        # loi de Mariotte : P x V = P1 x V1
    print(V, 'mL :', round(P), 'hPa')
\end{python}
\begin{console}
60 mL : 1013 hPa
50 mL : 1216 hPa
40 mL : 1520 hPa
30 mL : 2026 hPa
20 mL : 3039 hPa
\end{console}
\noindent\capacite{Capacité numérique} Le programme applique la loi de Mariotte : diviser le volume par 3 multiplie la pression par 3.
\essSection{4}{Loi fondamentale de la statique des fluides}
\begin{loi}[Statique des fluides]
Dans un fluide \textbf{incompressible} (un liquide) au repos, de masse volumique $\rho$, pour deux points d'altitudes $z_1$ et $z_2$
repérées sur un axe vertical \textbf{orienté vers le haut} :
\[P_2-P_1=\rho\,g\,(z_1-z_2).\]
À la profondeur $h$ sous la surface libre, $P=P_\mathrm{atm}+\rho\,g\,h$ : dans l'eau, la pression augmente de \qty{9,81e4}{Pa}, soit
environ \qty{1}{bar}, tous les \qty{10}{m}.
\end{loi}
\begin{propriete}[Conséquences]
La pression est la même en tous les points d'un même plan horizontal d'un liquide au repos. La surface libre d'un liquide
est horizontale ; dans des vases communicants, elle est au même niveau dans toutes les branches.
\end{propriete}
\par\penalty0
\schema{12-statique}
\begin{propriete}[Quelques ordres de grandeur]
\par\smallskip\noindent\hfill\begin{tabular}{lr@{\qquad}lr}
\toprule
Situation & Pression & Situation & Pression \\
\midrule
Sommet du mont Blanc & \qty{550}{hPa} & Autocuiseur en marche & \qty{2}{bar} \\
Niveau de la mer & \qty{1013}{hPa} & Plongeur à \qty{20}{m} & \qty{3}{bar} \\
Écart sur un étage (air) & \qty{0,35}{hPa} & Bouteille de plongée & \qty{200}{bar} \\
\bottomrule
\end{tabular}\hfill\null\par\smallskip
\end{propriete}
\begin{attention}
Dans $F=P\times S$, convertir les surfaces en \unit{m^2} ($\qty{1}{cm^2}=\qty{1e-4}{m^2}$). Dans la loi de la statique, bien vérifier
l'orientation de l'axe : vers le haut, le point le plus bas a la plus grande pression. La loi de Mariotte ne s'applique qu'aux
gaz, à température et quantité de matière constantes ; la loi de la statique ne s'applique qu'aux liquides.
\end{attention}
\begin{experience}[Tester les deux lois]
\textbf{Mariotte} : une seringue est reliée à un capteur de pression piloté par un microcontrôleur ; on relève la pression pour
différents volumes, en déplaçant lentement le piston pour que la température reste constante. \textbf{Statique} : on
enfonce la sonde d'un manomètre dans une éprouvette d'eau et l'on trace $P$ en fonction de la profondeur : on obtient
une droite de pente $\rho\,g$.
\end{experience}
\begin{rediger}
Écrire la loi littérale, préciser le système (quantité de gaz enfermée, liquide considéré) et les conditions d'application
(température constante, liquide au repos), convertir dans les unités du Système international, puis calculer.
\end{rediger}
\exercices

\rubrique{Découverte}
\begin{blocLivre}{ex-12-1}
\begin{exercice}[1]{Du microscopique au macroscopique}
Un automobiliste contrôle la pression de ses pneus le matin, à \qty{10}{\degreeCelsius}. Après une heure de route au soleil, l'air des pneus est à
\qty{40}{\degreeCelsius} ; le volume des pneus n'a pratiquement pas changé.
\begin{donnees}
$\rho_\mathrm{air}=\qty{1,2}{kg.m^{-3}}$ ; $\rho_\mathrm{eau}=\qty{1,00e3}{kg.m^{-3}}$.
\end{donnees}
\begin{questions}
\item Comment la température de l'air se traduit-elle à l'échelle des molécules ? Faire un schéma comparatif.
\item En déduire comment évolue la pression dans les pneus. Justifier à l'aide des chocs sur les parois.
\item Calculer la masse d'air contenue dans une pièce de \qty{5}{m} sur \qty{4}{m} et de \qty{2,5}{m} de hauteur.
\item L'eau liquide est environ 830 fois plus dense que l'air. Interpréter à l'échelle microscopique.
\end{questions}
\end{exercice}
\end{blocLivre}
\begin{blocLivre}{ex-12-2}
\begin{exercice}[1]{La force de l'air sur une fenêtre}
Une vitre rectangulaire mesure \qty{1,20}{m} sur \qty{0,80}{m}.
\begin{donnees}
$P_\mathrm{atm}=\qty{1,013e5}{Pa}$ ; $g=\qty{9,81}{N.kg^{-1}}$.
\end{donnees}
\begin{questions}
\item Calculer la valeur de la force pressante exercée par l'air extérieur sur la vitre.
\item Quelle masse aurait un objet dont le poids a cette valeur ?
\item Pourquoi la vitre ne se brise-t-elle pas ? Représenter les forces sur un schéma.
\item Pendant une tempête, la pression extérieure devient inférieure de \qty{30}{hPa} à la pression intérieure. Calculer la force
résultante sur la vitre et préciser son sens.
\end{questions}
\end{exercice}
\end{blocLivre}
\begin{blocLivre}{ex-12-3}
\begin{exercice}[1]{La seringue bouchée}
\schema{12-ex03}
\begin{donnees}
$P_\mathrm{atm}=\qty{1013}{hPa}$.
\end{donnees}
\begin{questions}
\item Quelle loi peut-on appliquer à l'air enfermé ? À quelles conditions ?
\item Calculer la pression de l'air enfermé lorsque son volume vaut \qty{20}{mL}.
\item Faire l'inventaire des forces qui s'exercent sur le piston et calculer la force que doit exercer l'opérateur.
\item Pour quel volume la pression de l'air enfermé vaut-elle \qty{2,0e5}{Pa} ?
\end{questions}
\end{exercice}
\end{blocLivre}
\begin{blocLivre}{ex-12-4}
\begin{exercice}[1]{La pression sous l'eau}
\begin{donnees}
$P_\mathrm{atm}=\qty{1,013e5}{Pa}$ ; $g=\qty{9,81}{N.kg^{-1}}$ ; eau douce : $\rho=\qty{1,00e3}{kg.m^{-3}}$ ; eau de mer : $\rho=\qty{1,03e3}{kg.m^{-3}}$.
\end{donnees}
\begin{questions}
\item Calculer la pression au fond d'une piscine de \qty{2,0}{m} de profondeur.
\item Calculer la pression à \qty{10}{m} de profondeur dans la mer. Comparer à la pression atmosphérique.
\item À quelle profondeur dans la mer la pression atteint-elle \qty{3,0e5}{Pa} ?
\item Tracer l'allure de la pression en fonction de la profondeur. De quel type de fonction s'agit-il ?
\end{questions}
\end{exercice}
\end{blocLivre}
\begin{blocLivre}{ex-12-5}
\begin{exercice}[1]{Des vases communicants}
\schema{12-ex05}
\begin{donnees}
$P_\mathrm{atm}=\qty{1,013e5}{Pa}$ ; $g=\qty{9,81}{N.kg^{-1}}$ ; $\rho_\mathrm{eau}=\qty{1,00e3}{kg.m^{-3}}$.
\end{donnees}
\begin{questions}
\item Classer par ordre croissant les pressions aux points A, C, D et E. Justifier.
\item Comparer les pressions en A et en B.
\item Calculer la pression en A puis en C.
\item Expliquer pourquoi la surface libre de l'eau est à la même hauteur dans les deux branches.
\end{questions}
\end{exercice}
\end{blocLivre}
\rubrique{Entraînement}
\begin{blocLivre}{ex-12-6}
\begin{exercice}[2]{La pression en plongée, en Python}
\begin{python}
rho, g, P0 = 1030, 9.81, 1.013e5     # SI : eau de mer
for h in range(0, 50, 10):            # profondeur en m
    P = P0 + rho * g * h
    print(h, 'm :', round(P / 1e5, 2), 'bar')
\end{python}
\begin{questions}
\item \capacite{Capacité numérique}Quelle relation le programme utilise-t-il ? Que représente \code{P0} ?
\item Prévoir les cinq lignes affichées.
\item De combien la pression augmente-t-elle tous les \qty{10}{m} ? À quelle profondeur est-elle le triple de la pression atmosphérique ?
\item Modifier le programme pour un lac de montagne d'eau douce où $P_0=\qty{0,80e5}{Pa}$. Que vaut la pression à \qty{20}{m} ?
\end{questions}
\end{exercice}
\end{blocLivre}
\begin{blocLivre}{ex-12-7}
\begin{exercice}[2]{Les poumons de l'apnéiste}
Un apnéiste inspire \qty{6,0}{L} d'air en surface, à la pression $P_0=\qty{1,013e5}{Pa}$, puis descend en bloquant sa respiration. On
admet que la température de l'air de ses poumons reste constante.
\begin{donnees}
$g=\qty{9,81}{N.kg^{-1}}$ ; eau de mer : $\rho=\qty{1,03e3}{kg.m^{-3}}$.
\end{donnees}
\begin{questions}
\item Calculer la pression à \qty{20}{m} de profondeur.
\item Calculer le volume de l'air de ses poumons à cette profondeur.
\item Tracer l'allure du volume des poumons en fonction de la profondeur, entre 0 et \qty{40}{m}.
\item Un plongeur équipé d'une bouteille respire à \qty{20}{m} de l'air à la pression ambiante. Pourquoi ne doit-il jamais
remonter en bloquant sa respiration ?
\end{questions}
\end{exercice}
\end{blocLivre}
\begin{blocLivre}{ex-12-8}
\begin{exercice}[2]{Le baromètre de Torricelli}
\schema{12-ex08}
\begin{donnees}
$\rho_\mathrm{mercure}=\qty{13,6e3}{kg.m^{-3}}$ ; $\rho_\mathrm{eau}=\qty{1,00e3}{kg.m^{-3}}$ ; $g=\qty{9,81}{N.kg^{-1}}$ ; $P_\mathrm{atm}=\qty{1,013e5}{Pa}$.
\end{donnees}
\begin{questions}
\item Que vaut la pression au sommet C de la colonne ? Et au point A du tube situé au niveau de la surface de la cuve ?
\item Appliquer la loi de la statique entre A et C et en déduire la pression atmosphérique.
\item Pourquoi parle-t-on encore de « millimètres de mercure » pour mesurer la pression ?
\item Quelle hauteur atteindrait la colonne avec de l'eau ? Pourquoi a-t-on choisi le mercure ?
\end{questions}
\end{exercice}
\end{blocLivre}
\begin{blocLivre}{ex-12-9}
\begin{exercice}[2]{La pression et l'altitude}
On admet, pour de petites hauteurs, que la masse volumique de l'air reste égale à \qty{1,2}{kg.m^{-3}}.
\begin{donnees}
$g=\qty{9,81}{N.kg^{-1}}$ ; pression au niveau de la mer : \qty{1013}{hPa}.
\end{donnees}
\begin{questions}
\item Calculer la diminution de pression quand on s'élève de \qty{100}{m}, puis de \qty{300}{m} (hauteur de la tour Eiffel).
\item Un téléphone muni d'un baromètre détecte des écarts de \qty{0,1}{hPa}. Peut-il repérer le passage d'un étage, haut de \qty{3}{m} ?
\item Calculer la pression au sommet du mont Blanc (\qty{4807}{m}) avec la même hypothèse.
\item On y mesure en réalité environ \qty{550}{hPa}. Quelle hypothèse est en défaut ? Justifier.
\end{questions}
\end{exercice}
\end{blocLivre}
\begin{blocLivre}{ex-12-10}
\begin{exercice}[2]{Mesurer une masse volumique avec un tube en U}
\schema{12-ex10}
\begin{donnees}
$\rho_\mathrm{eau}=\qty{1,00e3}{kg.m^{-3}}$ ; $g=\qty{9,81}{N.kg^{-1}}$ ; $P_\mathrm{atm}=\qty{1,013e5}{Pa}$.
\end{donnees}
\begin{questions}
\item On note A le point de l'interface eau-huile et B le point de l'eau situé au même niveau dans l'autre branche. Pourquoi
$P_\mathrm{A}=P_\mathrm{B}$ ?
\item Exprimer $P_\mathrm{A}$ en fonction de $P_\mathrm{atm}$, $\rho_\mathrm{h}$, $g$ et $h_\mathrm{h}$, puis $P_\mathrm{B}$ en fonction de $P_\mathrm{atm}$, $\rho_\mathrm{eau}$, $g$ et $h_\mathrm{e}$.
\item En déduire la masse volumique de l'huile.
\item Calculer la pression au point A.
\end{questions}
\end{exercice}
\end{blocLivre}
\begin{blocLivre}{ex-12-11}
\begin{exercice}[2]{Tester la loi de Mariotte}
\schema{12-ex11}
\par\smallskip\noindent\hfill\begin{tabular}{lcccccc}
\toprule
$V$ (mL) & 60 & 50 & 40 & 30 & 24 & 20 \\
$P$ (hPa) & 1012 & 1218 & 1515 & 2031 & 2526 & 3046 \\
\bottomrule
\end{tabular}\hfill\null\par\smallskip
\begin{questions}
\item Calculer le produit $P\times V$ pour chaque mesure. La loi de Mariotte est-elle vérifiée ?
\item Tracer $P$ en fonction de $\dfrac{1}{V}$. Quelle courbe obtient-on ?
\item Prévoir la pression pour un volume de \qty{15}{mL}.
\item Pourquoi faut-il déplacer le piston lentement ? Quel défaut du montage peut fausser les mesures ?
\end{questions}
\end{exercice}
\end{blocLivre}
\rubrique{Objectif DS}
\begin{blocLivre}{ex-12-12}
\begin{exercice}[3]{La plongée en bouteille}
\begin{docu}{1}{La bouteille}
Une bouteille de plongée de \qty{12}{L} est gonflée d'air à \qty{200}{bar}. Par sécurité, le plongeur doit remonter quand la pression
dans la bouteille descend à \qty{50}{bar} : c'est la réserve.
\end{docu}
\begin{docu}{2}{Le détendeur}
Le détendeur fournit au plongeur de l'air à la pression de l'eau qui l'entoure. Un plongeur au repos respire \qty{20}{L} d'air par
minute, quelle que soit la profondeur : ce volume est mesuré à la pression ambiante.
\end{docu}
\begin{donnees}
On prend $P_\mathrm{atm}=\qty{1,0}{bar}$ ; $g=\qty{9,81}{N.kg^{-1}}$ ; eau de mer : $\rho=\qty{1,03e3}{kg.m^{-3}}$ ; température constante.
\end{donnees}
\begin{questions}
\item Calculer la pression à \qty{20}{m} de profondeur, en bar.
\item Quel volume occuperait l'air de la bouteille à la pression atmosphérique ? Et l'air utilisable avant la réserve ?
\item À \qty{20}{m}, quel volume d'air, ramené à la pression atmosphérique, le plongeur consomme-t-il par minute ?
\item \initiative Calculer la durée de la plongée à \qty{20}{m} avant d'atteindre la réserve, et la comparer à celle d'une plongée près de la
surface.
\end{questions}
\end{exercice}
\end{blocLivre}
\begin{blocLivre}{ex-12-13}
\begin{exercice}[3]{Le ballon-sonde}
\begin{docu}{1}{Un ballon météorologique}
Chaque jour, des ballons gonflés à l'hélium emportent des instruments de mesure jusqu'à plus de \qty{25}{km} d'altitude. Au
lâcher, l'enveloppe de caoutchouc, presque sphérique, contient \qty{1,5}{m^3} d'hélium à \qty{1013}{hPa}.
\end{docu}
\begin{docu}{2}{L'altitude}
La pression de l'air vaut environ \qty{30}{hPa} à \qty{24}{km} d'altitude. Le ballon éclate quand son diamètre atteint \qty{6,0}{m} ; un parachute
ramène alors les instruments au sol.
\end{docu}
\begin{donnees}
On admet que l'hélium du ballon reste à la pression de l'air extérieur et que sa température est constante. Volume d'une
sphère : $V=\dfrac{4}{3}\,\pi\,r^3$.
\end{donnees}
\begin{questions}
\item Calculer le volume du ballon à \qty{24}{km} d'altitude. Combien de fois a-t-il augmenté ?
\item Calculer le diamètre du ballon au sol et à \qty{24}{km}.
\item \initiative Calculer la pression de l'air au moment où le ballon éclate.
\item En réalité, la température vaut environ \qty{-55}{\degreeCelsius} vers \qty{24}{km}. Critiquer le modèle utilisé.
\end{questions}
\end{exercice}
\end{blocLivre}
\begin{blocLivre}{ex-12-14}
\begin{exercice}[3]{Le château d'eau}
\begin{docu}{1}{Distribuer l'eau}
Un château d'eau stocke l'eau en hauteur : la pression qu'elle exerce dans les canalisations suffit à la distribuer sans pompe.
Pour un bon confort, la surpression au robinet (pression de l'eau moins pression atmosphérique) doit atteindre au moins
\qty{3,0}{bar}.
\end{docu}
\begin{docu}{2}{Le réseau}
\schema{12-ex14}
\end{docu}
\begin{donnees}
$\rho_\mathrm{eau}=\qty{1,00e3}{kg.m^{-3}}$ ; $g=\qty{9,81}{N.kg^{-1}}$ ; $\qty{1}{bar}=\qty{1e5}{Pa}$.
\end{donnees}
\begin{questions}
\item Calculer la surpression au robinet A de la maison.
\item Calculer la surpression au robinet B du dernier étage de l'immeuble. Le confort est-il assuré ?
\item \initiative À quelle altitude minimale devrait se trouver la surface de l'eau pour assurer le confort au robinet B ?
\item Une trappe de visite de \qty{0,50}{m^2} est placée au fond de la cuve. Calculer la force pressante supplémentaire due à l'eau.
\end{questions}
\end{exercice}
\end{blocLivre}
\begin{blocLivre}{ex-12-15}
\begin{exercice}[3]{Pourquoi les barrages sont-ils si épais à la base ?}
\begin{docu}{1}{Le barrage-poids}
\schema{12-ex15}
Un barrage-poids résiste à la poussée de l'eau grâce à sa masse : sa base, très large, peut dépasser \qty{40}{m} d'épaisseur.
\end{docu}
\begin{donnees}
$\rho_\mathrm{eau}=\qty{1,00e3}{kg.m^{-3}}$ ; $g=\qty{9,81}{N.kg^{-1}}$. On ne considère que la surpression due à l'eau, l'air agissant des deux côtés.
\end{donnees}
\begin{questions}
\item Calculer la surpression due à l'eau à mi-hauteur, puis au pied du barrage.
\item Représenter l'allure des forces pressantes sur la paroi verticale, de la surface au pied.
\item Calculer la force pressante sur une vanne de \qty{1,0}{m^2} placée au pied du barrage.
\item \initiative En utilisant la pression moyenne sur la paroi, estimer la force exercée par l'eau sur une bande verticale de paroi
de \qty{1}{m} de large. Conclure sur la forme du barrage.
\end{questions}
\end{exercice}
\end{blocLivre}
